\documentclass[11pt,a4paper]{article}
\usepackage[margin=19mm]{geometry}
\usepackage[colorlinks=true,urlcolor=blue]{hyperref}
\usepackage{parskip}
\newcommand{\record}[1]{\textbf{#1}\par\vspace{0.55\baselineskip}\hrule\vspace{0.35\baselineskip}}
\begin{document}

{\small IMAGING REFERRAL AND REPORTING WORKBOOK \hfill 1 of 16}

\section*{Plain wrist X-ray}

\subsection*{Before the scan}

\textbf{Clinical question}\par Is a wrist bone broken after the fall, and has it shifted?\par

\textbf{What to check first}\par Record the mechanism, time since injury, focal bone tenderness, swelling, deformity, finger movement, sensation and circulation.\par

\textbf{When to proceed or defer}\par When the assessment suggests fracture, wrist radiographs are usually the first imaging. If clinical suspicion is low, document the basis for deferral and the reassessment plan.\par

\textbf{When to defer or adapt the test}\par Use a child-sized technique for children. In pregnancy, review the indication and body region. A needed extremity X-ray can proceed when justified.\par

\record{Patient ID, date and ordering clinician}

\record{History received: exact account, source and relationship, recipient, date and time}

\record{Clinician assessment and prior tests, with dates and results}

\record{Guideline and exact scenario matched}

\record{Finding that supports ordering or deferral}

\record{Patient-specific reason for an exception, if applicable}

\record{Order or defer, urgency, and what the result would change}

\record{Safety decision and plain-language explanation for the patient}

\textbf{Basis for request}\par Supported clinical use, uncertain clinical benefit, research, or reason awaiting clarification. For research, record the study reference and applicable approvals.\par

\small\href{https://www.radiologyinfo.org/en/info/acs-acute-hand-and-wrist-trauma}{Clinical reference for this scenario}

\quad\href{https://pubmed.ncbi.nlm.nih.gov/28937575/}{Pregnancy and lactation}

\quad\href{https://www.fda.gov/radiation-emitting-products/medical-imaging/pediatric-x-ray-imaging}{Pediatric X-ray imaging}

\normalsize

\newpage

{\small IMAGING REFERRAL AND REPORTING WORKBOOK \hfill 1 of 16}

\section*{Plain wrist X-ray}

\subsection*{The scan-specific report}

\textbf{Clinical question}\par Is a wrist bone broken after the fall, and has it shifted?\par

\textbf{What this report should record}\par State side and views, adequacy, fracture site, joint involvement, tilt, displacement and shortening.\par

\record{Patient ID, date and reporting clinician}

\record{Exam performed, technique, side or site, and comparisons}

\record{Findings and relevant measurements}

\record{Image quality, coverage or other limitations}

\record{Impression that answers the clinical question}

\record{Question status: answered, partly answered or inconclusive}

\record{Follow-up action, responsible clinician or team, and due date}

\record{Urgent communication: recipient, time and method}

\small\href{https://www.radiologyinfo.org/en/info/acs-acute-hand-and-wrist-trauma}{Clinical reference for this scenario}

\normalsize

\newpage

{\small IMAGING REFERRAL AND REPORTING WORKBOOK \hfill 1 of 16}

\section*{Plain wrist X-ray}

\subsection*{What the image shows and means}

\textbf{Clinical question}\par Is a wrist bone broken after the fall, and has it shifted?\par

\textbf{Interrupted bone cortex}\par Visible fracture\par

\textbf{Fragment tilt in degrees or shift in millimetres}\par Altered alignment\par

\textbf{Fracture line reaching the joint}\par Joint surface involvement\par

\textbf{What the image leaves unanswered}\par If radiographs are normal but suspicion persists, repeat radiographs, CT or MRI may be appropriate.\par

\textbf{Illustration guide for a 4:3 plate}\par Show lateral examples with and without a break. Label the radius and mark the tilt with guide lines.\par

\textbf{Worked interpretation example}\par Finding: Interrupted bone cortex. Interpretation: visible fracture. If this finding is absent, use the limitations above to decide whether the original question remains open.\par

\textbf{Other plain X-ray examples}\par Chest: assess an opacity when pneumonia is suspected. Joint: assess joint-space narrowing in arthritis. Tube check: identify the tube tip and assess its position.\par

\record{Finding, location, measurement and unit}

\record{Does the finding support the hypothesis, make it less likely, or leave it unresolved?}

\record{Plain-language explanation}

\record{Audit challenge: reviewer and date, finding disputed, evidence, response or resolution}

\small\href{https://www.radiologyinfo.org/en/info/acs-acute-hand-and-wrist-trauma}{Clinical reference for this scenario}

\normalsize

\newpage

{\small IMAGING REFERRAL AND REPORTING WORKBOOK \hfill 1 of 16}

\section*{Plain wrist X-ray}

\subsection*{Find similar cases and check the evidence}

\textbf{Clinical question}\par Is a wrist bone broken after the fall, and has it shifted?\par

\textbf{Optional review aid}\par Use de-identified clinical terms. Keep names, identifiers, reports and patient images out of public searches. Record examples for review without delaying care.\par

\textbf{1. Check the decision using only pre-scan information}\par fall on outstretched hand focal wrist tenderness initial radiography guideline\par

\textbf{2. Find a teaching case using the imaging findings}\par wrist radiograph distal radius fracture dorsal angulation teaching case\par

\record{Actual search terms, source link and access date}

\record{What matches and what differs}

\record{How the diagnosis was established, or whether confirmation was unclear}

\record{What the source adds and what remains uncertain}

\small\href{https://www.radiologyinfo.org/en/info/acs-acute-hand-and-wrist-trauma}{Clinical reference for this scenario}

\normalsize

\newpage

{\small IMAGING REFERRAL AND REPORTING WORKBOOK \hfill 2 of 16}

\section*{Head CT}

\subsection*{Before the scan}

\textbf{Clinical question}\par Is there bleeding inside the skull after a mild head injury?\par

\textbf{What to check first}\par Record mechanism, alertness, neurological findings, vomiting or amnesia, and relevant medicines. Identify whether a validated decision rule applies.\par

\textbf{When to proceed or defer}\par Use noncontrast CT when indicated by the applicable rule or urgent clinical findings. A qualifying low-risk assessment can support deferral.\par

\textbf{When to defer or adapt the test}\par Use a pediatric decision rule and size-adjusted technique for children. In pregnancy, justify exposure with the imaging team. For contrast CT, review kidney function when indicated and any prior reaction to iodinated contrast.\par

\record{Patient ID, date and ordering clinician}

\record{History received: exact account, source and relationship, recipient, date and time}

\record{Clinician assessment and prior tests, with dates and results}

\record{Guideline and exact scenario matched}

\record{Finding that supports ordering or deferral}

\record{Patient-specific reason for an exception, if applicable}

\record{Order or defer, urgency, and what the result would change}

\record{Safety decision and plain-language explanation for the patient}

\textbf{Basis for request}\par Supported clinical use, uncertain clinical benefit, research, or reason awaiting clarification. For research, record the study reference and applicable approvals.\par

\small\href{https://acsearch.acr.org/docs/69481/Narrative}{Clinical reference for this scenario}

\quad\href{https://pubmed.ncbi.nlm.nih.gov/28937575/}{Pregnancy and lactation}

\quad\href{https://www.fda.gov/radiation-emitting-products/medical-imaging/pediatric-x-ray-imaging}{Pediatric X-ray imaging}

\quad\href{https://www.acr.org/clinical-resources/clinical-tools-and-reference/contrast-manual}{Contrast safety}

\normalsize

\newpage

{\small IMAGING REFERRAL AND REPORTING WORKBOOK \hfill 2 of 16}

\section*{Head CT}

\subsection*{The scan-specific report}

\textbf{Clinical question}\par Is there bleeding inside the skull after a mild head injury?\par

\textbf{What this report should record}\par Record protocol, image quality, bleeding location and type, mass effect, midline shift and skull injury if seen.\par

\record{Patient ID, date and reporting clinician}

\record{Exam performed, technique, side or site, and comparisons}

\record{Findings and relevant measurements}

\record{Image quality, coverage or other limitations}

\record{Impression that answers the clinical question}

\record{Question status: answered, partly answered or inconclusive}

\record{Follow-up action, responsible clinician or team, and due date}

\record{Urgent communication: recipient, time and method}

\small\href{https://acsearch.acr.org/docs/69481/Narrative}{Clinical reference for this scenario}

\normalsize

\newpage

{\small IMAGING REFERRAL AND REPORTING WORKBOOK \hfill 2 of 16}

\section*{Head CT}

\subsection*{What the image shows and means}

\textbf{Clinical question}\par Is there bleeding inside the skull after a mild head injury?\par

\textbf{Bright collection outside the brain}\par May indicate acute extra-axial blood\par

\textbf{Bright focus in brain tissue}\par May indicate contusion or haemorrhage\par

\textbf{Midline displacement in millimetres}\par Mass effect\par

\textbf{What the image leaves unanswered}\par A normal CT can coexist with concussion. Document any limitation and answer the acute bleeding question explicitly.\par

\textbf{Illustration guide for a 4:3 plate}\par Label an illustrative axial slice and use an arrow to identify the collection and a line to mark the midline.\par

\textbf{Worked interpretation example}\par Finding: Bright collection outside the brain. Interpretation: may indicate acute extra-axial blood. If this finding is absent, use the limitations above to decide whether the original question remains open.\par

\record{Finding, location, measurement and unit}

\record{Does the finding support the hypothesis, make it less likely, or leave it unresolved?}

\record{Plain-language explanation}

\record{Audit challenge: reviewer and date, finding disputed, evidence, response or resolution}

\small\href{https://acsearch.acr.org/docs/69481/Narrative}{Clinical reference for this scenario}

\normalsize

\newpage

{\small IMAGING REFERRAL AND REPORTING WORKBOOK \hfill 2 of 16}

\section*{Head CT}

\subsection*{Find similar cases and check the evidence}

\textbf{Clinical question}\par Is there bleeding inside the skull after a mild head injury?\par

\textbf{Optional review aid}\par Use de-identified clinical terms. Keep names, identifiers, reports and patient images out of public searches. Record examples for review without delaying care.\par

\textbf{1. Check the decision using only pre-scan information}\par adult mild head injury validated decision rule noncontrast CT guideline\par

\textbf{2. Find a teaching case using the imaging findings}\par noncontrast head CT acute subdural hematoma midline shift teaching case\par

\record{Actual search terms, source link and access date}

\record{What matches and what differs}

\record{How the diagnosis was established, or whether confirmation was unclear}

\record{What the source adds and what remains uncertain}

\small\href{https://acsearch.acr.org/docs/69481/Narrative}{Clinical reference for this scenario}

\normalsize

\newpage

{\small IMAGING REFERRAL AND REPORTING WORKBOOK \hfill 3 of 16}

\section*{Knee MRI}

\subsection*{Before the scan}

\textbf{Clinical question}\par What is causing persistent knee pain after initial evaluation?\par

\textbf{What to check first}\par Examine the knee, document the location of pain and functional limitation, and review plain radiographs.\par

\textbf{When to proceed or defer}\par If radiographs answer the management question, consider deferral. If they are negative or show an effusion and the cause remains unclear, MRI without contrast is usually appropriate in that scenario.\par

\textbf{When to defer or adapt the test}\par Verify implants and suspected metallic foreign bodies before entering the scanner. Follow device-specific conditions or choose an alternative if safety cannot be established. For gadolinium, review pregnancy, kidney disease and prior reactions.\par

\record{Patient ID, date and ordering clinician}

\record{History received: exact account, source and relationship, recipient, date and time}

\record{Clinician assessment and prior tests, with dates and results}

\record{Guideline and exact scenario matched}

\record{Finding that supports ordering or deferral}

\record{Patient-specific reason for an exception, if applicable}

\record{Order or defer, urgency, and what the result would change}

\record{Safety decision and plain-language explanation for the patient}

\textbf{Basis for request}\par Supported clinical use, uncertain clinical benefit, research, or reason awaiting clarification. For research, record the study reference and applicable approvals.\par

\small\href{https://acsearch.acr.org/docs/69432/Narrative}{Clinical reference for this scenario}

\quad\href{https://www.acr.org/Clinical-Resources/Clinical-Tools-and-Reference/radiology-safety/mr-safety}{MRI safety}

\quad\href{https://www.acr.org/clinical-resources/clinical-tools-and-reference/contrast-manual}{Contrast safety}

\quad\href{https://pubmed.ncbi.nlm.nih.gov/28937575/}{Pregnancy and lactation}

\normalsize

\newpage

{\small IMAGING REFERRAL AND REPORTING WORKBOOK \hfill 3 of 16}

\section*{Knee MRI}

\subsection*{The scan-specific report}

\textbf{Clinical question}\par What is causing persistent knee pain after initial evaluation?\par

\textbf{What this report should record}\par Record side, sequences, menisci, ligaments, cartilage, marrow and joint fluid. Identify pertinent abnormalities and limitations.\par

\record{Patient ID, date and reporting clinician}

\record{Exam performed, technique, side or site, and comparisons}

\record{Findings and relevant measurements}

\record{Image quality, coverage or other limitations}

\record{Impression that answers the clinical question}

\record{Question status: answered, partly answered or inconclusive}

\record{Follow-up action, responsible clinician or team, and due date}

\record{Urgent communication: recipient, time and method}

\small\href{https://acsearch.acr.org/docs/69432/Narrative}{Clinical reference for this scenario}

\normalsize

\newpage

{\small IMAGING REFERRAL AND REPORTING WORKBOOK \hfill 3 of 16}

\section*{Knee MRI}

\subsection*{What the image shows and means}

\textbf{Clinical question}\par What is causing persistent knee pain after initial evaluation?\par

\textbf{Meniscal signal reaching an articular surface}\par Supports a tear\par

\textbf{Marrow signal change}\par Can reflect injury or other causes\par

\textbf{Cartilage loss location and extent}\par Helps localise joint damage\par

\textbf{What the image leaves unanswered}\par Imaging abnormalities can be incidental. Relate findings to the location and pattern of symptoms.\par

\textbf{Illustration guide for a 4:3 plate}\par Show a labelled sagittal schematic with an inset marking a meniscal tear and a separate area of cartilage loss.\par

\textbf{Worked interpretation example}\par Finding: Meniscal signal reaching an articular surface. Interpretation: supports a tear. If this finding is absent, use the limitations above to decide whether the original question remains open.\par

\record{Finding, location, measurement and unit}

\record{Does the finding support the hypothesis, make it less likely, or leave it unresolved?}

\record{Plain-language explanation}

\record{Audit challenge: reviewer and date, finding disputed, evidence, response or resolution}

\small\href{https://acsearch.acr.org/docs/69432/Narrative}{Clinical reference for this scenario}

\normalsize

\newpage

{\small IMAGING REFERRAL AND REPORTING WORKBOOK \hfill 3 of 16}

\section*{Knee MRI}

\subsection*{Find similar cases and check the evidence}

\textbf{Clinical question}\par What is causing persistent knee pain after initial evaluation?\par

\textbf{Optional review aid}\par Use de-identified clinical terms. Keep names, identifiers, reports and patient images out of public searches. Record examples for review without delaying care.\par

\textbf{1. Check the decision using only pre-scan information}\par adult chronic knee pain negative radiographs MRI appropriateness guideline\par

\textbf{2. Find a teaching case using the imaging findings}\par knee MRI meniscal tear articular surface teaching case\par

\record{Actual search terms, source link and access date}

\record{What matches and what differs}

\record{How the diagnosis was established, or whether confirmation was unclear}

\record{What the source adds and what remains uncertain}

\small\href{https://acsearch.acr.org/docs/69432/Narrative}{Clinical reference for this scenario}

\normalsize

\newpage

{\small IMAGING REFERRAL AND REPORTING WORKBOOK \hfill 4 of 16}

\section*{Swallowing fluoroscopy}

\subsection*{Before the scan}

\textbf{Clinical question}\par Does food or liquid enter the airway during swallowing?\par

\textbf{What to check first}\par Identify whether the difficulty is initiating a swallow or a sensation of food sticking lower down. Review a clinical swallowing assessment when available.\par

\textbf{When to proceed or defer}\par A modified barium swallow is appropriate for selected oropharyngeal swallowing problems. Choose the examination that matches the suspected site and question.\par

\textbf{When to defer or adapt the test}\par Review pregnancy and tailor radiation exposure for children. If the patient is unstable or cannot safely participate, review timing and swallowing assessment options with the team.\par

\record{Patient ID, date and ordering clinician}

\record{History received: exact account, source and relationship, recipient, date and time}

\record{Clinician assessment and prior tests, with dates and results}

\record{Guideline and exact scenario matched}

\record{Finding that supports ordering or deferral}

\record{Patient-specific reason for an exception, if applicable}

\record{Order or defer, urgency, and what the result would change}

\record{Safety decision and plain-language explanation for the patient}

\textbf{Basis for request}\par Supported clinical use, uncertain clinical benefit, research, or reason awaiting clarification. For research, record the study reference and applicable approvals.\par

\small\href{https://acsearch.acr.org/docs/69471/Narrative}{Clinical reference for this scenario}

\quad\href{https://pubmed.ncbi.nlm.nih.gov/28937575/}{Pregnancy and lactation}

\quad\href{https://www.fda.gov/radiation-emitting-products/medical-imaging/pediatric-x-ray-imaging}{Pediatric X-ray imaging}

\normalsize

\newpage

{\small IMAGING REFERRAL AND REPORTING WORKBOOK \hfill 4 of 16}

\section*{Swallowing fluoroscopy}

\subsection*{The scan-specific report}

\textbf{Clinical question}\par Does food or liquid enter the airway during swallowing?\par

\textbf{What this report should record}\par Record consistencies, views, airway entry, residue, timing and effect of tested strategies.\par

\record{Patient ID, date and reporting clinician}

\record{Exam performed, technique, side or site, and comparisons}

\record{Findings and relevant measurements}

\record{Image quality, coverage or other limitations}

\record{Impression that answers the clinical question}

\record{Question status: answered, partly answered or inconclusive}

\record{Follow-up action, responsible clinician or team, and due date}

\record{Urgent communication: recipient, time and method}

\small\href{https://acsearch.acr.org/docs/69471/Narrative}{Clinical reference for this scenario}

\normalsize

\newpage

{\small IMAGING REFERRAL AND REPORTING WORKBOOK \hfill 4 of 16}

\section*{Swallowing fluoroscopy}

\subsection*{What the image shows and means}

\textbf{Clinical question}\par Does food or liquid enter the airway during swallowing?\par

\textbf{Contrast below vocal folds}\par Aspiration\par

\textbf{Contrast entering airway above vocal folds}\par Penetration\par

\textbf{Residue after swallow}\par Incomplete clearance\par

\textbf{What the image leaves unanswered}\par The study samples particular swallows. State which consistencies or strategies could be assessed.\par

\textbf{Illustration guide for a 4:3 plate}\par Use three frames showing bolus passage and label the airway, vocal folds and oesophagus.\par

\textbf{Worked interpretation example}\par Finding: Contrast below vocal folds. Interpretation: aspiration. If this finding is absent, use the limitations above to decide whether the original question remains open.\par

\record{Finding, location, measurement and unit}

\record{Does the finding support the hypothesis, make it less likely, or leave it unresolved?}

\record{Plain-language explanation}

\record{Audit challenge: reviewer and date, finding disputed, evidence, response or resolution}

\small\href{https://acsearch.acr.org/docs/69471/Narrative}{Clinical reference for this scenario}

\normalsize

\newpage

{\small IMAGING REFERRAL AND REPORTING WORKBOOK \hfill 4 of 16}

\section*{Swallowing fluoroscopy}

\subsection*{Find similar cases and check the evidence}

\textbf{Clinical question}\par Does food or liquid enter the airway during swallowing?\par

\textbf{Optional review aid}\par Use de-identified clinical terms. Keep names, identifiers, reports and patient images out of public searches. Record examples for review without delaying care.\par

\textbf{1. Check the decision using only pre-scan information}\par suspected oropharyngeal dysphagia modified barium swallow appropriateness guideline\par

\textbf{2. Find a teaching case using the imaging findings}\par videofluoroscopic swallowing study contrast below vocal folds aspiration teaching case\par

\record{Actual search terms, source link and access date}

\record{What matches and what differs}

\record{How the diagnosis was established, or whether confirmation was unclear}

\record{What the source adds and what remains uncertain}

\small\href{https://acsearch.acr.org/docs/69471/Narrative}{Clinical reference for this scenario}

\normalsize

\newpage

{\small IMAGING REFERRAL AND REPORTING WORKBOOK \hfill 5 of 16}

\section*{Leg vein ultrasound}

\subsection*{Before the scan}

\textbf{Clinical question}\par Is there a deep vein clot in a symptomatic leg?\par

\textbf{What to check first}\par In the NICE two-level pathway, score clot likelihood. Use D-dimer first in the unlikely group.\par

\textbf{When to proceed or defer}\par Negative D-dimer in the unlikely group generally avoids ultrasound. Positive D-dimer or a likely score leads to ultrasound via the applicable pathway.\par

\textbf{When to defer or adapt the test}\par Adapt positioning for frailty, pain or wounds. Use a pregnancy-specific or pediatric clot pathway when applicable. If views are incomplete, document the limitation and decide how to complete the assessment.\par

\record{Patient ID, date and ordering clinician}

\record{History received: exact account, source and relationship, recipient, date and time}

\record{Clinician assessment and prior tests, with dates and results}

\record{Guideline and exact scenario matched}

\record{Finding that supports ordering or deferral}

\record{Patient-specific reason for an exception, if applicable}

\record{Order or defer, urgency, and what the result would change}

\record{Safety decision and plain-language explanation for the patient}

\textbf{Basis for request}\par Supported clinical use, uncertain clinical benefit, research, or reason awaiting clarification. For research, record the study reference and applicable approvals.\par

\small\href{https://www.nice.org.uk/guidance/ng158/chapter/Recommendations}{Clinical reference for this scenario}

\normalsize

\newpage

{\small IMAGING REFERRAL AND REPORTING WORKBOOK \hfill 5 of 16}

\section*{Leg vein ultrasound}

\subsection*{The scan-specific report}

\textbf{Clinical question}\par Is there a deep vein clot in a symptomatic leg?\par

\textbf{What this report should record}\par List side, venous segments examined, compressibility, Doppler findings and technical limitations. Specify whether the protocol examined proximal veins or the whole leg.\par

\record{Patient ID, date and reporting clinician}

\record{Exam performed, technique, side or site, and comparisons}

\record{Findings and relevant measurements}

\record{Image quality, coverage or other limitations}

\record{Impression that answers the clinical question}

\record{Question status: answered, partly answered or inconclusive}

\record{Follow-up action, responsible clinician or team, and due date}

\record{Urgent communication: recipient, time and method}

\small\href{https://www.nice.org.uk/guidance/ng158/chapter/Recommendations}{Clinical reference for this scenario}

\normalsize

\newpage

{\small IMAGING REFERRAL AND REPORTING WORKBOOK \hfill 5 of 16}

\section*{Leg vein ultrasound}

\subsection*{What the image shows and means}

\textbf{Clinical question}\par Is there a deep vein clot in a symptomatic leg?\par

\textbf{Vein fails to compress}\par Key sign supporting thrombosis\par

\textbf{Visible intraluminal material}\par Helps localise thrombus\par

\textbf{Abnormal flow}\par May support obstruction in context\par

\textbf{What the image leaves unanswered}\par The meaning of a negative result depends on scan coverage and pretest probability. Some negative proximal scans require repeat imaging under the pathway.\par

\textbf{Illustration guide for a 4:3 plate}\par Show paired compression views of a patent and a thrombosed vein.\par

\textbf{Worked interpretation example}\par Finding: Vein fails to compress. Interpretation: key sign supporting thrombosis. If this finding is absent, use the limitations above to decide whether the original question remains open.\par

\record{Finding, location, measurement and unit}

\record{Does the finding support the hypothesis, make it less likely, or leave it unresolved?}

\record{Plain-language explanation}

\record{Audit challenge: reviewer and date, finding disputed, evidence, response or resolution}

\small\href{https://www.nice.org.uk/guidance/ng158/chapter/Recommendations}{Clinical reference for this scenario}

\normalsize

\newpage

{\small IMAGING REFERRAL AND REPORTING WORKBOOK \hfill 5 of 16}

\section*{Leg vein ultrasound}

\subsection*{Find similar cases and check the evidence}

\textbf{Clinical question}\par Is there a deep vein clot in a symptomatic leg?\par

\textbf{Optional review aid}\par Use de-identified clinical terms. Keep names, identifiers, reports and patient images out of public searches. Record examples for review without delaying care.\par

\textbf{1. Check the decision using only pre-scan information}\par suspected leg DVT Wells D-dimer ultrasound referral guideline\par

\textbf{2. Find a teaching case using the imaging findings}\par compression ultrasound noncompressible femoral vein deep vein thrombosis teaching case\par

\record{Actual search terms, source link and access date}

\record{What matches and what differs}

\record{How the diagnosis was established, or whether confirmation was unclear}

\record{What the source adds and what remains uncertain}

\small\href{https://www.nice.org.uk/guidance/ng158/chapter/Recommendations}{Clinical reference for this scenario}

\normalsize

\newpage

{\small IMAGING REFERRAL AND REPORTING WORKBOOK \hfill 6 of 16}

\section*{Retinal OCT}

\subsection*{Before the scan}

\textbf{Clinical question}\par Is there fluid or thickening in the macula?\par

\textbf{What to check first}\par Review vision symptoms, acuity, retinal examination and previous OCT when available.\par

\textbf{When to proceed or defer}\par Use OCT when suspected macular swelling or change would affect management. Review recent adequate images before ordering a duplicate.\par

\textbf{When to defer or adapt the test}\par Adapt fixation and positioning for children, frailty or limited mobility. If image quality is inadequate, repeat with adaptations or record which measurements cannot be interpreted.\par

\record{Patient ID, date and ordering clinician}

\record{History received: exact account, source and relationship, recipient, date and time}

\record{Clinician assessment and prior tests, with dates and results}

\record{Guideline and exact scenario matched}

\record{Finding that supports ordering or deferral}

\record{Patient-specific reason for an exception, if applicable}

\record{Order or defer, urgency, and what the result would change}

\record{Safety decision and plain-language explanation for the patient}

\textbf{Basis for request}\par Supported clinical use, uncertain clinical benefit, research, or reason awaiting clarification. For research, record the study reference and applicable approvals.\par

\small\href{https://www.aaojournal.org/article/S0161-6420(19)32902-5/fulltext/}{Clinical reference for this scenario}

\normalsize

\newpage

{\small IMAGING REFERRAL AND REPORTING WORKBOOK \hfill 6 of 16}

\section*{Retinal OCT}

\subsection*{The scan-specific report}

\textbf{Clinical question}\par Is there fluid or thickening in the macula?\par

\textbf{What this report should record}\par Record eye, scan quality, macular thickness when reliable, fluid location, and change from baseline.\par

\record{Patient ID, date and reporting clinician}

\record{Exam performed, technique, side or site, and comparisons}

\record{Findings and relevant measurements}

\record{Image quality, coverage or other limitations}

\record{Impression that answers the clinical question}

\record{Question status: answered, partly answered or inconclusive}

\record{Follow-up action, responsible clinician or team, and due date}

\record{Urgent communication: recipient, time and method}

\small\href{https://www.aaojournal.org/article/S0161-6420(19)32902-5/fulltext/}{Clinical reference for this scenario}

\normalsize

\newpage

{\small IMAGING REFERRAL AND REPORTING WORKBOOK \hfill 6 of 16}

\section*{Retinal OCT}

\subsection*{What the image shows and means}

\textbf{Clinical question}\par Is there fluid or thickening in the macula?\par

\textbf{Dark pockets within the retina}\par Intraretinal fluid\par

\textbf{Dark space beneath the retina}\par Subretinal fluid\par

\textbf{Central thickness in micrometres}\par Tracking measure when comparable\par

\textbf{What the image leaves unanswered}\par Thickness alone does not establish the cause. Motion and segmentation errors can affect measurements.\par

\textbf{Illustration guide for a 4:3 plate}\par Label retinal layers and outline a fluid pocket on a comparison schematic.\par

\textbf{Worked interpretation example}\par Finding: Dark pockets within the retina. Interpretation: intraretinal fluid. If this finding is absent, use the limitations above to decide whether the original question remains open.\par

\record{Finding, location, measurement and unit}

\record{Does the finding support the hypothesis, make it less likely, or leave it unresolved?}

\record{Plain-language explanation}

\record{Audit challenge: reviewer and date, finding disputed, evidence, response or resolution}

\small\href{https://www.aaojournal.org/article/S0161-6420(19)32902-5/fulltext/}{Clinical reference for this scenario}

\normalsize

\newpage

{\small IMAGING REFERRAL AND REPORTING WORKBOOK \hfill 6 of 16}

\section*{Retinal OCT}

\subsection*{Find similar cases and check the evidence}

\textbf{Clinical question}\par Is there fluid or thickening in the macula?\par

\textbf{Optional review aid}\par Use de-identified clinical terms. Keep names, identifiers, reports and patient images out of public searches. Record examples for review without delaying care.\par

\textbf{1. Check the decision using only pre-scan information}\par suspected diabetic macular edema OCT evaluation guideline\par

\textbf{2. Find a teaching case using the imaging findings}\par retinal OCT diabetic macular edema intraretinal fluid teaching case\par

\record{Actual search terms, source link and access date}

\record{What matches and what differs}

\record{How the diagnosis was established, or whether confirmation was unclear}

\record{What the source adds and what remains uncertain}

\small\href{https://www.aaojournal.org/article/S0161-6420(19)32902-5/fulltext/}{Clinical reference for this scenario}

\normalsize

\newpage

{\small IMAGING REFERRAL AND REPORTING WORKBOOK \hfill 7 of 16}

\section*{FDG PET CT}

\subsection*{Before the scan}

\textbf{Clinical question}\par Where is biopsy-confirmed Hodgkin lymphoma present before treatment?\par

\textbf{What to check first}\par Confirm pathology and the staging question. Review relevant prior imaging and record factors affecting tracer uptake or interpretation.\par

\textbf{When to proceed or defer}\par Use FDG PET/CT for initial staging of confirmed Hodgkin lymphoma. Lack of symptoms does not remove the staging indication.\par

\textbf{When to defer or adapt the test}\par Review pregnancy and breastfeeding with nuclear medicine before tracer administration. Tailor administered activity for children. Check glucose and preparation against the local FDG protocol.\par

\record{Patient ID, date and ordering clinician}

\record{History received: exact account, source and relationship, recipient, date and time}

\record{Clinician assessment and prior tests, with dates and results}

\record{Guideline and exact scenario matched}

\record{Finding that supports ordering or deferral}

\record{Patient-specific reason for an exception, if applicable}

\record{Order or defer, urgency, and what the result would change}

\record{Safety decision and plain-language explanation for the patient}

\textbf{Basis for request}\par Supported clinical use, uncertain clinical benefit, research, or reason awaiting clarification. For research, record the study reference and applicable approvals.\par

\small\href{https://pmc.ncbi.nlm.nih.gov/articles/12165998/}{Clinical reference for this scenario}

\quad\href{https://pubmed.ncbi.nlm.nih.gov/28937575/}{Pregnancy and lactation}

\normalsize

\newpage

{\small IMAGING REFERRAL AND REPORTING WORKBOOK \hfill 7 of 16}

\section*{FDG PET CT}

\subsection*{The scan-specific report}

\textbf{Clinical question}\par Where is biopsy-confirmed Hodgkin lymphoma present before treatment?\par

\textbf{What this report should record}\par Record tracer, scan coverage, relevant node groups and organ sites, CT correlates and limits. Distinguish suspicious from indeterminate uptake.\par

\record{Patient ID, date and reporting clinician}

\record{Exam performed, technique, side or site, and comparisons}

\record{Findings and relevant measurements}

\record{Image quality, coverage or other limitations}

\record{Impression that answers the clinical question}

\record{Question status: answered, partly answered or inconclusive}

\record{Follow-up action, responsible clinician or team, and due date}

\record{Urgent communication: recipient, time and method}

\small\href{https://pmc.ncbi.nlm.nih.gov/articles/12165998/}{Clinical reference for this scenario}

\normalsize

\newpage

{\small IMAGING REFERRAL AND REPORTING WORKBOOK \hfill 7 of 16}

\section*{FDG PET CT}

\subsection*{What the image shows and means}

\textbf{Clinical question}\par Where is biopsy-confirmed Hodgkin lymphoma present before treatment?\par

\textbf{Focal FDG uptake mapped to a node}\par May indicate lymphoma involvement\par

\textbf{Nodal sites above and below the diaphragm}\par Contributes to staging\par

\textbf{Uptake in an organ}\par Requires pattern and anatomic correlation\par

\textbf{What the image leaves unanswered}\par Infection and inflammation can also increase tracer uptake. Correlate atypical findings with the clinical context.\par

\textbf{Illustration guide for a 4:3 plate}\par Show paired PET and fused PET/CT schematics with the same nodal site marked.\par

\textbf{Worked interpretation example}\par Finding: Focal FDG uptake mapped to a node. Interpretation: may indicate lymphoma involvement. If this finding is absent, use the limitations above to decide whether the original question remains open.\par

\record{Finding, location, measurement and unit}

\record{Does the finding support the hypothesis, make it less likely, or leave it unresolved?}

\record{Plain-language explanation}

\record{Audit challenge: reviewer and date, finding disputed, evidence, response or resolution}

\small\href{https://pmc.ncbi.nlm.nih.gov/articles/12165998/}{Clinical reference for this scenario}

\normalsize

\newpage

{\small IMAGING REFERRAL AND REPORTING WORKBOOK \hfill 7 of 16}

\section*{FDG PET CT}

\subsection*{Find similar cases and check the evidence}

\textbf{Clinical question}\par Where is biopsy-confirmed Hodgkin lymphoma present before treatment?\par

\textbf{Optional review aid}\par Use de-identified clinical terms. Keep names, identifiers, reports and patient images out of public searches. Record examples for review without delaying care.\par

\textbf{1. Check the decision using only pre-scan information}\par biopsy confirmed Hodgkin lymphoma initial staging FDG PET CT guideline\par

\textbf{2. Find a teaching case using the imaging findings}\par Hodgkin lymphoma baseline FDG PET CT nodal involvement teaching case\par

\record{Actual search terms, source link and access date}

\record{What matches and what differs}

\record{How the diagnosis was established, or whether confirmation was unclear}

\record{What the source adds and what remains uncertain}

\small\href{https://pmc.ncbi.nlm.nih.gov/articles/12165998/}{Clinical reference for this scenario}

\normalsize

\newpage

{\small IMAGING REFERRAL AND REPORTING WORKBOOK \hfill 8 of 16}

\section*{Cardiac stress SPECT}

\subsection*{Before the scan}

\textbf{Clinical question}\par Is there reduced blood flow to part of the heart muscle on stress?\par

\textbf{What to check first}\par For stable symptoms, estimate coronary disease likelihood and review resting ECG, exercise capacity and known coronary disease.\par

\textbf{When to proceed or defer}\par Low-risk stable presentations may defer testing. When stress imaging is indicated, document why SPECT fits. Acute symptoms follow a different pathway.\par

\textbf{When to defer or adapt the test}\par Review pregnancy and breastfeeding before tracer administration. Assess contraindications to the chosen exercise or pharmacologic stress method. Active or unstable symptoms require a separate clinical pathway.\par

\record{Patient ID, date and ordering clinician}

\record{History received: exact account, source and relationship, recipient, date and time}

\record{Clinician assessment and prior tests, with dates and results}

\record{Guideline and exact scenario matched}

\record{Finding that supports ordering or deferral}

\record{Patient-specific reason for an exception, if applicable}

\record{Order or defer, urgency, and what the result would change}

\record{Safety decision and plain-language explanation for the patient}

\textbf{Basis for request}\par Supported clinical use, uncertain clinical benefit, research, or reason awaiting clarification. For research, record the study reference and applicable approvals.\par

\small\href{https://www.jacc.org/doi/10.1016/j.jacc.2021.07.053}{Clinical reference for this scenario}

\quad\href{https://pubmed.ncbi.nlm.nih.gov/28937575/}{Pregnancy and lactation}

\normalsize

\newpage

{\small IMAGING REFERRAL AND REPORTING WORKBOOK \hfill 8 of 16}

\section*{Cardiac stress SPECT}

\subsection*{The scan-specific report}

\textbf{Clinical question}\par Is there reduced blood flow to part of the heart muscle on stress?\par

\textbf{What this report should record}\par Record stress method and adequacy, tracer, location and extent of defects, reversibility, ventricular function if measured, and artefacts.\par

\record{Patient ID, date and reporting clinician}

\record{Exam performed, technique, side or site, and comparisons}

\record{Findings and relevant measurements}

\record{Image quality, coverage or other limitations}

\record{Impression that answers the clinical question}

\record{Question status: answered, partly answered or inconclusive}

\record{Follow-up action, responsible clinician or team, and due date}

\record{Urgent communication: recipient, time and method}

\small\href{https://www.jacc.org/doi/10.1016/j.jacc.2021.07.053}{Clinical reference for this scenario}

\normalsize

\newpage

{\small IMAGING REFERRAL AND REPORTING WORKBOOK \hfill 8 of 16}

\section*{Cardiac stress SPECT}

\subsection*{What the image shows and means}

\textbf{Clinical question}\par Is there reduced blood flow to part of the heart muscle on stress?\par

\textbf{Stress defect improving at rest}\par Supports inducible ischaemia\par

\textbf{Defect on both studies}\par May reflect scar or artefact\par

\textbf{Ejection fraction}\par Describes pumping function when measured\par

\textbf{What the image leaves unanswered}\par Attenuation and motion can mimic defects. A normal study does not exclude every cause of chest symptoms.\par

\textbf{Illustration guide for a 4:3 plate}\par Place corresponding stress and rest maps side by side with a matched colour scale.\par

\textbf{Worked interpretation example}\par Finding: Stress defect improving at rest. Interpretation: supports inducible ischaemia. If this finding is absent, use the limitations above to decide whether the original question remains open.\par

\record{Finding, location, measurement and unit}

\record{Does the finding support the hypothesis, make it less likely, or leave it unresolved?}

\record{Plain-language explanation}

\record{Audit challenge: reviewer and date, finding disputed, evidence, response or resolution}

\small\href{https://www.jacc.org/doi/10.1016/j.jacc.2021.07.053}{Clinical reference for this scenario}

\normalsize

\newpage

{\small IMAGING REFERRAL AND REPORTING WORKBOOK \hfill 8 of 16}

\section*{Cardiac stress SPECT}

\subsection*{Find similar cases and check the evidence}

\textbf{Clinical question}\par Is there reduced blood flow to part of the heart muscle on stress?\par

\textbf{Optional review aid}\par Use de-identified clinical terms. Keep names, identifiers, reports and patient images out of public searches. Record examples for review without delaying care.\par

\textbf{1. Check the decision using only pre-scan information}\par stable chest pain intermediate risk stress myocardial perfusion imaging guideline\par

\textbf{2. Find a teaching case using the imaging findings}\par stress rest SPECT reversible perfusion defect teaching case\par

\record{Actual search terms, source link and access date}

\record{What matches and what differs}

\record{How the diagnosis was established, or whether confirmation was unclear}

\record{What the source adds and what remains uncertain}

\small\href{https://www.jacc.org/doi/10.1016/j.jacc.2021.07.053}{Clinical reference for this scenario}

\normalsize

\newpage

{\small IMAGING REFERRAL AND REPORTING WORKBOOK \hfill 9 of 16}

\section*{Mammography and tomosynthesis}

\subsection*{Before the scan}

\textbf{Clinical question}\par Is screening due, and is there a finding that requires recall?\par

\textbf{What to check first}\par Confirm screening eligibility, interval, risk and prior studies. Ask about focal symptoms and check whether a diagnostic pathway applies.\par

\textbf{When to proceed or defer}\par Screen when due under the named guideline even in an asymptomatic patient. Use the diagnostic pathway for a new focal breast concern.\par

\textbf{When to defer or adapt the test}\par Review new symptoms and pregnancy or breastfeeding with breast imaging. Routine screening eligibility and urgent diagnostic assessment are separate decisions.\par

\record{Patient ID, date and ordering clinician}

\record{History received: exact account, source and relationship, recipient, date and time}

\record{Clinician assessment and prior tests, with dates and results}

\record{Guideline and exact scenario matched}

\record{Finding that supports ordering or deferral}

\record{Patient-specific reason for an exception, if applicable}

\record{Order or defer, urgency, and what the result would change}

\record{Safety decision and plain-language explanation for the patient}

\textbf{Basis for request}\par Supported clinical use, uncertain clinical benefit, research, or reason awaiting clarification. For research, record the study reference and applicable approvals.\par

\small\href{https://www.uspreventiveservicestaskforce.org/uspstf/recommendation/breast-cancer-screening}{Clinical reference for this scenario}

\quad\href{https://pubmed.ncbi.nlm.nih.gov/28937575/}{Pregnancy and lactation}

\normalsize

\newpage

{\small IMAGING REFERRAL AND REPORTING WORKBOOK \hfill 9 of 16}

\section*{Mammography and tomosynthesis}

\subsection*{The scan-specific report}

\textbf{Clinical question}\par Is screening due, and is there a finding that requires recall?\par

\textbf{What this report should record}\par Record technique, comparison, breast density, lesion descriptors, BI-RADS assessment and next action.\par

\record{Patient ID, date and reporting clinician}

\record{Exam performed, technique, side or site, and comparisons}

\record{Findings and relevant measurements}

\record{Image quality, coverage or other limitations}

\record{Impression that answers the clinical question}

\record{Question status: answered, partly answered or inconclusive}

\record{Follow-up action, responsible clinician or team, and due date}

\record{Urgent communication: recipient, time and method}

\small\href{https://www.uspreventiveservicestaskforce.org/uspstf/recommendation/breast-cancer-screening}{Clinical reference for this scenario}

\normalsize

\newpage

{\small IMAGING REFERRAL AND REPORTING WORKBOOK \hfill 9 of 16}

\section*{Mammography and tomosynthesis}

\subsection*{What the image shows and means}

\textbf{Clinical question}\par Is screening due, and is there a finding that requires recall?\par

\textbf{Grouped calcifications}\par Morphology and distribution guide assessment\par

\textbf{Mass shape and margins}\par Inform concern\par

\textbf{Architectural distortion}\par May warrant further evaluation\par

\textbf{What the image leaves unanswered}\par A screening recall is a request for further assessment. Persistent focal symptoms need clinical follow-up even after negative screening.\par

\textbf{Illustration guide for a 4:3 plate}\par Show schematic examples of a mass, calcifications and distortion with separate labels.\par

\textbf{Worked interpretation example}\par Finding: Grouped calcifications. Interpretation: morphology and distribution guide assessment. If this finding is absent, use the limitations above to decide whether the original question remains open.\par

\record{Finding, location, measurement and unit}

\record{Does the finding support the hypothesis, make it less likely, or leave it unresolved?}

\record{Plain-language explanation}

\record{Audit challenge: reviewer and date, finding disputed, evidence, response or resolution}

\small\href{https://www.uspreventiveservicestaskforce.org/uspstf/recommendation/breast-cancer-screening}{Clinical reference for this scenario}

\normalsize

\newpage

{\small IMAGING REFERRAL AND REPORTING WORKBOOK \hfill 9 of 16}

\section*{Mammography and tomosynthesis}

\subsection*{Find similar cases and check the evidence}

\textbf{Clinical question}\par Is screening due, and is there a finding that requires recall?\par

\textbf{Optional review aid}\par Use de-identified clinical terms. Keep names, identifiers, reports and patient images out of public searches. Record examples for review without delaying care.\par

\textbf{1. Check the decision using only pre-scan information}\par asymptomatic average risk breast screening mammography tomosynthesis guideline\par

\textbf{2. Find a teaching case using the imaging findings}\par mammography grouped calcifications morphology teaching case\par

\record{Actual search terms, source link and access date}

\record{What matches and what differs}

\record{How the diagnosis was established, or whether confirmation was unclear}

\record{What the source adds and what remains uncertain}

\small\href{https://www.uspreventiveservicestaskforce.org/uspstf/recommendation/breast-cancer-screening}{Clinical reference for this scenario}

\normalsize

\newpage

{\small IMAGING REFERRAL AND REPORTING WORKBOOK \hfill 10 of 16}

\section*{DXA bone density}

\subsection*{Before the scan}

\textbf{Clinical question}\par Would measuring bone density inform osteoporosis assessment?\par

\textbf{What to check first}\par Review age, fracture history, medicines, risk factors and previous DXA. Identify the screening criteria that apply.\par

\textbf{When to proceed or defer}\par Proceed when the applicable screening criteria are met or the diagnostic result could change care. For repeat scans, document timing and expected benefit.\par

\textbf{When to defer or adapt the test}\par Check pregnancy before elective DXA and review whether it should be deferred. In children, use pediatric reporting standards. Consider positioning, prior implants and artefacts in older adults.\par

\record{Patient ID, date and ordering clinician}

\record{History received: exact account, source and relationship, recipient, date and time}

\record{Clinician assessment and prior tests, with dates and results}

\record{Guideline and exact scenario matched}

\record{Finding that supports ordering or deferral}

\record{Patient-specific reason for an exception, if applicable}

\record{Order or defer, urgency, and what the result would change}

\record{Safety decision and plain-language explanation for the patient}

\textbf{Basis for request}\par Supported clinical use, uncertain clinical benefit, research, or reason awaiting clarification. For research, record the study reference and applicable approvals.\par

\small\href{https://www.uspreventiveservicestaskforce.org/uspstf/recommendation/osteoporosis-screening}{Clinical reference for this scenario}

\quad\href{https://pubmed.ncbi.nlm.nih.gov/28937575/}{Pregnancy and lactation}

\normalsize

\newpage

{\small IMAGING REFERRAL AND REPORTING WORKBOOK \hfill 10 of 16}

\section*{DXA bone density}

\subsection*{The scan-specific report}

\textbf{Clinical question}\par Would measuring bone density inform osteoporosis assessment?\par

\textbf{What this report should record}\par Record measured sites, bone mineral density, applicable scores and limitations. Compare only with a technically suitable prior study.\par

\record{Patient ID, date and reporting clinician}

\record{Exam performed, technique, side or site, and comparisons}

\record{Findings and relevant measurements}

\record{Image quality, coverage or other limitations}

\record{Impression that answers the clinical question}

\record{Question status: answered, partly answered or inconclusive}

\record{Follow-up action, responsible clinician or team, and due date}

\record{Urgent communication: recipient, time and method}

\small\href{https://www.uspreventiveservicestaskforce.org/uspstf/recommendation/osteoporosis-screening}{Clinical reference for this scenario}

\normalsize

\newpage

{\small IMAGING REFERRAL AND REPORTING WORKBOOK \hfill 10 of 16}

\section*{DXA bone density}

\subsection*{What the image shows and means}

\textbf{Clinical question}\par Would measuring bone density inform osteoporosis assessment?\par

\textbf{Bone mineral density in g/cm$^2$}\par Mineral per projected area\par

\textbf{T-score at a valid site}\par Supports classification in applicable adults\par

\textbf{Change beyond measurement variability}\par May indicate a real interval change\par

\textbf{What the image leaves unanswered}\par Positioning, arthritis and other artefacts may distort values. Fracture risk also depends on factors beyond density.\par

\textbf{Illustration guide for a 4:3 plate}\par Show hip and spine regions of interest and label where the measurement is taken.\par

\textbf{Worked interpretation example}\par Finding: Bone mineral density in g/cm$^2$. Interpretation: mineral per projected area. If this finding is absent, use the limitations above to decide whether the original question remains open.\par

\record{Finding, location, measurement and unit}

\record{Does the finding support the hypothesis, make it less likely, or leave it unresolved?}

\record{Plain-language explanation}

\record{Audit challenge: reviewer and date, finding disputed, evidence, response or resolution}

\small\href{https://www.uspreventiveservicestaskforce.org/uspstf/recommendation/osteoporosis-screening}{Clinical reference for this scenario}

\normalsize

\newpage

{\small IMAGING REFERRAL AND REPORTING WORKBOOK \hfill 10 of 16}

\section*{DXA bone density}

\subsection*{Find similar cases and check the evidence}

\textbf{Clinical question}\par Would measuring bone density inform osteoporosis assessment?\par

\textbf{Optional review aid}\par Use de-identified clinical terms. Keep names, identifiers, reports and patient images out of public searches. Record examples for review without delaying care.\par

\textbf{1. Check the decision using only pre-scan information}\par osteoporosis risk screening eligibility DXA guideline\par

\textbf{2. Find a teaching case using the imaging findings}\par hip spine DXA T-score positioning artifact teaching case\par

\record{Actual search terms, source link and access date}

\record{What matches and what differs}

\record{How the diagnosis was established, or whether confirmation was unclear}

\record{What the source adds and what remains uncertain}

\small\href{https://www.uspreventiveservicestaskforce.org/uspstf/recommendation/osteoporosis-screening}{Clinical reference for this scenario}

\normalsize

\newpage

{\small IMAGING REFERRAL AND REPORTING WORKBOOK \hfill 11 of 16}

\section*{HIDA gallbladder scan}

\subsection*{Before the scan}

\textbf{Clinical question}\par Does persistent suspicion after ultrasound warrant testing gallbladder filling?\par

\textbf{What to check first}\par Review symptoms, exam, ultrasound and relevant labs. Record preparation and medicines that can influence tracer transit.\par

\textbf{When to proceed or defer}\par Ultrasound usually comes first. If it answers the question, defer HIDA. If acute cholecystitis remains suspected after equivocal ultrasound, HIDA is one option.\par

\textbf{When to defer or adapt the test}\par Review pregnancy, breastfeeding and child-specific activity with nuclear medicine. Check fasting duration and medicines that can alter gallbladder filling.\par

\record{Patient ID, date and ordering clinician}

\record{History received: exact account, source and relationship, recipient, date and time}

\record{Clinician assessment and prior tests, with dates and results}

\record{Guideline and exact scenario matched}

\record{Finding that supports ordering or deferral}

\record{Patient-specific reason for an exception, if applicable}

\record{Order or defer, urgency, and what the result would change}

\record{Safety decision and plain-language explanation for the patient}

\textbf{Basis for request}\par Supported clinical use, uncertain clinical benefit, research, or reason awaiting clarification. For research, record the study reference and applicable approvals.\par

\small\href{https://www.radiologyinfo.org/en/info/acs-right-upper-quadrant-pain}{Clinical reference for this scenario}

\quad\href{https://pubmed.ncbi.nlm.nih.gov/28937575/}{Pregnancy and lactation}

\normalsize

\newpage

{\small IMAGING REFERRAL AND REPORTING WORKBOOK \hfill 11 of 16}

\section*{HIDA gallbladder scan}

\subsection*{The scan-specific report}

\textbf{Clinical question}\par Does persistent suspicion after ultrasound warrant testing gallbladder filling?\par

\textbf{What this report should record}\par Record tracer, timing of liver uptake and biliary transit, gallbladder visualisation, delayed images or interventions and limitations.\par

\record{Patient ID, date and reporting clinician}

\record{Exam performed, technique, side or site, and comparisons}

\record{Findings and relevant measurements}

\record{Image quality, coverage or other limitations}

\record{Impression that answers the clinical question}

\record{Question status: answered, partly answered or inconclusive}

\record{Follow-up action, responsible clinician or team, and due date}

\record{Urgent communication: recipient, time and method}

\small\href{https://www.radiologyinfo.org/en/info/acs-right-upper-quadrant-pain}{Clinical reference for this scenario}

\normalsize

\newpage

{\small IMAGING REFERRAL AND REPORTING WORKBOOK \hfill 11 of 16}

\section*{HIDA gallbladder scan}

\subsection*{What the image shows and means}

\textbf{Clinical question}\par Does persistent suspicion after ultrasound warrant testing gallbladder filling?\par

\textbf{Gallbladder appears}\par Supports a patent cystic duct\par

\textbf{Persistent nonvisualisation under an adequate protocol}\par Supports cystic duct obstruction in the appropriate setting\par

\textbf{Bowel activity}\par Documents tracer transit into intestine\par

\textbf{What the image leaves unanswered}\par Fasting, medicines or severe illness can affect the study. Interpret nonvisualisation in clinical context.\par

\textbf{Illustration guide for a 4:3 plate}\par Draw a timed sequence of liver, ducts, gallbladder and bowel with clear labels.\par

\textbf{Worked interpretation example}\par Finding: Gallbladder appears. Interpretation: supports a patent cystic duct. If this finding is absent, use the limitations above to decide whether the original question remains open.\par

\record{Finding, location, measurement and unit}

\record{Does the finding support the hypothesis, make it less likely, or leave it unresolved?}

\record{Plain-language explanation}

\record{Audit challenge: reviewer and date, finding disputed, evidence, response or resolution}

\small\href{https://www.radiologyinfo.org/en/info/acs-right-upper-quadrant-pain}{Clinical reference for this scenario}

\normalsize

\newpage

{\small IMAGING REFERRAL AND REPORTING WORKBOOK \hfill 11 of 16}

\section*{HIDA gallbladder scan}

\subsection*{Find similar cases and check the evidence}

\textbf{Clinical question}\par Does persistent suspicion after ultrasound warrant testing gallbladder filling?\par

\textbf{Optional review aid}\par Use de-identified clinical terms. Keep names, identifiers, reports and patient images out of public searches. Record examples for review without delaying care.\par

\textbf{1. Check the decision using only pre-scan information}\par suspected acute cholecystitis equivocal ultrasound HIDA appropriateness guideline\par

\textbf{2. Find a teaching case using the imaging findings}\par HIDA gallbladder nonvisualization acute cholecystitis teaching case\par

\record{Actual search terms, source link and access date}

\record{What matches and what differs}

\record{How the diagnosis was established, or whether confirmation was unclear}

\record{What the source adds and what remains uncertain}

\small\href{https://www.radiologyinfo.org/en/info/acs-right-upper-quadrant-pain}{Clinical reference for this scenario}

\normalsize

\newpage

{\small IMAGING REFERRAL AND REPORTING WORKBOOK \hfill 12 of 16}

\section*{Diuretic renal scintigraphy}

\subsection*{Before the scan}

\textbf{Clinical question}\par Is the dilated kidney obstructed, and how much function does it contribute?\par

\textbf{What to check first}\par Review ultrasound or other anatomic imaging, creatinine, prior surgery and drainage history.\par

\textbf{When to proceed or defer}\par Use when dilation leaves a functional question and measuring drainage or relative function could change care.\par

\textbf{When to defer or adapt the test}\par Review pregnancy, breastfeeding, hydration, kidney function and child-specific activity. Confirm whether the patient can void or needs an adapted drainage protocol.\par

\record{Patient ID, date and ordering clinician}

\record{History received: exact account, source and relationship, recipient, date and time}

\record{Clinician assessment and prior tests, with dates and results}

\record{Guideline and exact scenario matched}

\record{Finding that supports ordering or deferral}

\record{Patient-specific reason for an exception, if applicable}

\record{Order or defer, urgency, and what the result would change}

\record{Safety decision and plain-language explanation for the patient}

\textbf{Basis for request}\par Supported clinical use, uncertain clinical benefit, research, or reason awaiting clarification. For research, record the study reference and applicable approvals.\par

\small\href{https://snmmi.org/Web/Web/Clinical-Practice/Procedure-Standards/Standards/Diuretic-Renal-Scintigraphy-in-Adults-With-Suspected-Upper-Urinary-Tract-Obstruction-1-0.aspx}{Clinical reference for this scenario}

\quad\href{https://pubmed.ncbi.nlm.nih.gov/28937575/}{Pregnancy and lactation}

\normalsize

\newpage

{\small IMAGING REFERRAL AND REPORTING WORKBOOK \hfill 12 of 16}

\section*{Diuretic renal scintigraphy}

\subsection*{The scan-specific report}

\textbf{Clinical question}\par Is the dilated kidney obstructed, and how much function does it contribute?\par

\textbf{What this report should record}\par Record tracer, hydration, diuretic timing, relative function, curves, post-void images when relevant and overall interpretation.\par

\record{Patient ID, date and reporting clinician}

\record{Exam performed, technique, side or site, and comparisons}

\record{Findings and relevant measurements}

\record{Image quality, coverage or other limitations}

\record{Impression that answers the clinical question}

\record{Question status: answered, partly answered or inconclusive}

\record{Follow-up action, responsible clinician or team, and due date}

\record{Urgent communication: recipient, time and method}

\small\href{https://snmmi.org/Web/Web/Clinical-Practice/Procedure-Standards/Standards/Diuretic-Renal-Scintigraphy-in-Adults-With-Suspected-Upper-Urinary-Tract-Obstruction-1-0.aspx}{Clinical reference for this scenario}

\normalsize

\newpage

{\small IMAGING REFERRAL AND REPORTING WORKBOOK \hfill 12 of 16}

\section*{Diuretic renal scintigraphy}

\subsection*{What the image shows and means}

\textbf{Clinical question}\par Is the dilated kidney obstructed, and how much function does it contribute?\par

\textbf{Relative uptake percentage}\par Estimates each kidney's contribution\par

\textbf{Tracer washes out after diuretic}\par Supports drainage\par

\textbf{Retained tracer with supporting findings}\par May support impaired drainage\par

\textbf{What the image leaves unanswered}\par Slow washout alone is insufficient. Function, bladder fullness and technique can affect interpretation.\par

\textbf{Illustration guide for a 4:3 plate}\par Pair the images with labelled time-activity curves for contrasting drainage patterns.\par

\textbf{Worked interpretation example}\par Finding: Relative uptake percentage. Interpretation: estimates each kidney's contribution. If this finding is absent, use the limitations above to decide whether the original question remains open.\par

\record{Finding, location, measurement and unit}

\record{Does the finding support the hypothesis, make it less likely, or leave it unresolved?}

\record{Plain-language explanation}

\record{Audit challenge: reviewer and date, finding disputed, evidence, response or resolution}

\small\href{https://snmmi.org/Web/Web/Clinical-Practice/Procedure-Standards/Standards/Diuretic-Renal-Scintigraphy-in-Adults-With-Suspected-Upper-Urinary-Tract-Obstruction-1-0.aspx}{Clinical reference for this scenario}

\normalsize

\newpage

{\small IMAGING REFERRAL AND REPORTING WORKBOOK \hfill 12 of 16}

\section*{Diuretic renal scintigraphy}

\subsection*{Find similar cases and check the evidence}

\textbf{Clinical question}\par Is the dilated kidney obstructed, and how much function does it contribute?\par

\textbf{Optional review aid}\par Use de-identified clinical terms. Keep names, identifiers, reports and patient images out of public searches. Record examples for review without delaying care.\par

\textbf{1. Check the decision using only pre-scan information}\par hydronephrosis suspected urinary obstruction diuretic renal scintigraphy guideline\par

\textbf{2. Find a teaching case using the imaging findings}\par MAG3 diuretic renogram delayed washout teaching case\par

\record{Actual search terms, source link and access date}

\record{What matches and what differs}

\record{How the diagnosis was established, or whether confirmation was unclear}

\record{What the source adds and what remains uncertain}

\small\href{https://snmmi.org/Web/Web/Clinical-Practice/Procedure-Standards/Standards/Diuretic-Renal-Scintigraphy-in-Adults-With-Suspected-Upper-Urinary-Tract-Obstruction-1-0.aspx}{Clinical reference for this scenario}

\normalsize

\newpage

{\small IMAGING REFERRAL AND REPORTING WORKBOOK \hfill 13 of 16}

\section*{Upper endoscopy}

\subsection*{Before the scan}

\textbf{Clinical question}\par Does persistent indigestion warrant direct inspection of the upper digestive tract?\par

\textbf{What to check first}\par Review symptoms, age, cancer risk and earlier treatment. In many lower-risk adults under 60, noninvasive H. pylori testing precedes endoscopy under ACG/CAG guidance.\par

\textbf{When to proceed or defer}\par Initial testing and treatment can precede endoscopy in suitable patients. Age, risk and concerning findings may change the pathway.\par

\textbf{When to defer or adapt the test}\par Assess cardiopulmonary stability, aspiration risk, fasting, and sedation needs. In pregnancy or frailty, discuss timing and sedation options with the procedural team.\par

\record{Patient ID, date and ordering clinician}

\record{History received: exact account, source and relationship, recipient, date and time}

\record{Clinician assessment and prior tests, with dates and results}

\record{Guideline and exact scenario matched}

\record{Finding that supports ordering or deferral}

\record{Patient-specific reason for an exception, if applicable}

\record{Order or defer, urgency, and what the result would change}

\record{Safety decision and plain-language explanation for the patient}

\textbf{Basis for request}\par Supported clinical use, uncertain clinical benefit, research, or reason awaiting clarification. For research, record the study reference and applicable approvals.\par

\small\href{https://gi.org/guideline/management-of-dyspepsia}{Clinical reference for this scenario}

\quad\href{https://pmc.ncbi.nlm.nih.gov/articles/PMC10850688/}{Endoscopy sedation}

\normalsize

\newpage

{\small IMAGING REFERRAL AND REPORTING WORKBOOK \hfill 13 of 16}

\section*{Upper endoscopy}

\subsection*{The scan-specific report}

\textbf{Clinical question}\par Does persistent indigestion warrant direct inspection of the upper digestive tract?\par

\textbf{What this report should record}\par Record extent of examination, views and limitations, lesion site and size, biopsy sites and pending pathology.\par

\record{Patient ID, date and reporting clinician}

\record{Exam performed, technique, side or site, and comparisons}

\record{Findings and relevant measurements}

\record{Image quality, coverage or other limitations}

\record{Impression that answers the clinical question}

\record{Question status: answered, partly answered or inconclusive}

\record{Follow-up action, responsible clinician or team, and due date}

\record{Urgent communication: recipient, time and method}

\small\href{https://gi.org/guideline/management-of-dyspepsia}{Clinical reference for this scenario}

\normalsize

\newpage

{\small IMAGING REFERRAL AND REPORTING WORKBOOK \hfill 13 of 16}

\section*{Upper endoscopy}

\subsection*{What the image shows and means}

\textbf{Clinical question}\par Does persistent indigestion warrant direct inspection of the upper digestive tract?\par

\textbf{Mucosal crater}\par Visually consistent with an ulcer\par

\textbf{Bleeding or visible vessel}\par Characterises bleeding features\par

\textbf{Irregular lesion}\par May require tissue diagnosis\par

\textbf{What the image leaves unanswered}\par A visually normal lining can coexist with microscopic disease. Keep visual findings and pathology results distinct.\par

\textbf{Illustration guide for a 4:3 plate}\par Show a mucosal view and label an ulcer edge, base and biopsy location conceptually.\par

\textbf{Worked interpretation example}\par Finding: Mucosal crater. Interpretation: visually consistent with an ulcer. If this finding is absent, use the limitations above to decide whether the original question remains open.\par

\record{Finding, location, measurement and unit}

\record{Does the finding support the hypothesis, make it less likely, or leave it unresolved?}

\record{Plain-language explanation}

\record{Audit challenge: reviewer and date, finding disputed, evidence, response or resolution}

\small\href{https://gi.org/guideline/management-of-dyspepsia}{Clinical reference for this scenario}

\normalsize

\newpage

{\small IMAGING REFERRAL AND REPORTING WORKBOOK \hfill 13 of 16}

\section*{Upper endoscopy}

\subsection*{Find similar cases and check the evidence}

\textbf{Clinical question}\par Does persistent indigestion warrant direct inspection of the upper digestive tract?\par

\textbf{Optional review aid}\par Use de-identified clinical terms. Keep names, identifiers, reports and patient images out of public searches. Record examples for review without delaying care.\par

\textbf{1. Check the decision using only pre-scan information}\par adult persistent dyspepsia H pylori initial evaluation endoscopy guideline\par

\textbf{2. Find a teaching case using the imaging findings}\par upper endoscopy gastric ulcer mucosal crater teaching case\par

\record{Actual search terms, source link and access date}

\record{What matches and what differs}

\record{How the diagnosis was established, or whether confirmation was unclear}

\record{What the source adds and what remains uncertain}

\small\href{https://gi.org/guideline/management-of-dyspepsia}{Clinical reference for this scenario}

\normalsize

\newpage

{\small IMAGING REFERRAL AND REPORTING WORKBOOK \hfill 14 of 16}

\section*{Endoscopic ultrasound}

\subsection*{Before the scan}

\textbf{Clinical question}\par Does a solid pancreatic lesion require closer characterisation or tissue sampling?\par

\textbf{What to check first}\par Review dedicated CT or MRI when available. Specify what remains unresolved and whether tissue would affect management.\par

\textbf{When to proceed or defer}\par Proceed when EUS adds clinically relevant information or tissue acquisition is indicated. Review whether earlier tests have already answered the management question.\par

\textbf{When to defer or adapt the test}\par Assess aspiration risk, cardiopulmonary stability and sedation. If sampling is planned, review bleeding risk and antithrombotic medicines with the procedural team.\par

\record{Patient ID, date and ordering clinician}

\record{History received: exact account, source and relationship, recipient, date and time}

\record{Clinician assessment and prior tests, with dates and results}

\record{Guideline and exact scenario matched}

\record{Finding that supports ordering or deferral}

\record{Patient-specific reason for an exception, if applicable}

\record{Order or defer, urgency, and what the result would change}

\record{Safety decision and plain-language explanation for the patient}

\textbf{Basis for request}\par Supported clinical use, uncertain clinical benefit, research, or reason awaiting clarification. For research, record the study reference and applicable approvals.\par

\small\href{https://pubmed.ncbi.nlm.nih.gov/39387777/}{Clinical reference for this scenario}

\quad\href{https://pmc.ncbi.nlm.nih.gov/articles/PMC10850688/}{Endoscopy sedation}

\normalsize

\newpage

{\small IMAGING REFERRAL AND REPORTING WORKBOOK \hfill 14 of 16}

\section*{Endoscopic ultrasound}

\subsection*{The scan-specific report}

\textbf{Clinical question}\par Does a solid pancreatic lesion require closer characterisation or tissue sampling?\par

\textbf{What this report should record}\par Record lesion site and dimensions, visible duct and vessel relationships, sampling method and pending pathology.\par

\record{Patient ID, date and reporting clinician}

\record{Exam performed, technique, side or site, and comparisons}

\record{Findings and relevant measurements}

\record{Image quality, coverage or other limitations}

\record{Impression that answers the clinical question}

\record{Question status: answered, partly answered or inconclusive}

\record{Follow-up action, responsible clinician or team, and due date}

\record{Urgent communication: recipient, time and method}

\small\href{https://pubmed.ncbi.nlm.nih.gov/39387777/}{Clinical reference for this scenario}

\normalsize

\newpage

{\small IMAGING REFERRAL AND REPORTING WORKBOOK \hfill 14 of 16}

\section*{Endoscopic ultrasound}

\subsection*{What the image shows and means}

\textbf{Clinical question}\par Does a solid pancreatic lesion require closer characterisation or tissue sampling?\par

\textbf{Measured lesion dimensions}\par Characterises size\par

\textbf{Duct dilation or interruption}\par May reflect obstruction\par

\textbf{Sample result}\par Can establish tissue diagnosis when adequate\par

\textbf{What the image leaves unanswered}\par The ultrasound appearance alone may not identify the lesion type. An inadequate sample leaves the question open.\par

\textbf{Illustration guide for a 4:3 plate}\par Draw the probe beside the pancreas and show a magnified inset with the lesion and biopsy path.\par

\textbf{Worked interpretation example}\par Finding: Measured lesion dimensions. Interpretation: characterises size. If this finding is absent, use the limitations above to decide whether the original question remains open.\par

\record{Finding, location, measurement and unit}

\record{Does the finding support the hypothesis, make it less likely, or leave it unresolved?}

\record{Plain-language explanation}

\record{Audit challenge: reviewer and date, finding disputed, evidence, response or resolution}

\small\href{https://pubmed.ncbi.nlm.nih.gov/39387777/}{Clinical reference for this scenario}

\normalsize

\newpage

{\small IMAGING REFERRAL AND REPORTING WORKBOOK \hfill 14 of 16}

\section*{Endoscopic ultrasound}

\subsection*{Find similar cases and check the evidence}

\textbf{Clinical question}\par Does a solid pancreatic lesion require closer characterisation or tissue sampling?\par

\textbf{Optional review aid}\par Use de-identified clinical terms. Keep names, identifiers, reports and patient images out of public searches. Record examples for review without delaying care.\par

\textbf{1. Check the decision using only pre-scan information}\par solid pancreatic lesion on CT endoscopic ultrasound tissue acquisition guideline\par

\textbf{2. Find a teaching case using the imaging findings}\par endoscopic ultrasound solid pancreatic mass duct obstruction teaching case\par

\record{Actual search terms, source link and access date}

\record{What matches and what differs}

\record{How the diagnosis was established, or whether confirmation was unclear}

\record{What the source adds and what remains uncertain}

\small\href{https://pubmed.ncbi.nlm.nih.gov/39387777/}{Clinical reference for this scenario}

\normalsize

\newpage

{\small IMAGING REFERRAL AND REPORTING WORKBOOK \hfill 15 of 16}

\section*{Fundus photography}

\subsection*{Before the scan}

\textbf{Clinical question}\par Are there visible signs of diabetic retinopathy on screening?\par

\textbf{What to check first}\par Check diabetes history, previous screening, new symptoms and the screening programme's eligibility and interval.\par

\textbf{When to proceed or defer}\par Take photographs when screening is due in a photography-based programme. New symptoms may need a diagnostic eye assessment.\par

\textbf{When to defer or adapt the test}\par Adjust for child cooperation or limited mobility. If drops are planned, review dilation precautions, including suspected narrow angles, with the eye team.\par

\record{Patient ID, date and ordering clinician}

\record{History received: exact account, source and relationship, recipient, date and time}

\record{Clinician assessment and prior tests, with dates and results}

\record{Guideline and exact scenario matched}

\record{Finding that supports ordering or deferral}

\record{Patient-specific reason for an exception, if applicable}

\record{Order or defer, urgency, and what the result would change}

\record{Safety decision and plain-language explanation for the patient}

\textbf{Basis for request}\par Supported clinical use, uncertain clinical benefit, research, or reason awaiting clarification. For research, record the study reference and applicable approvals.\par

\small\href{https://www.aaojournal.org/article/S0161-6420(19)32902-5/fulltext/}{Clinical reference for this scenario}

\normalsize

\newpage

{\small IMAGING REFERRAL AND REPORTING WORKBOOK \hfill 15 of 16}

\section*{Fundus photography}

\subsection*{The scan-specific report}

\textbf{Clinical question}\par Are there visible signs of diabetic retinopathy on screening?\par

\textbf{What this report should record}\par Record eye, photographed fields, image quality, visible lesions, screening grade and recommended referral or recall.\par

\record{Patient ID, date and reporting clinician}

\record{Exam performed, technique, side or site, and comparisons}

\record{Findings and relevant measurements}

\record{Image quality, coverage or other limitations}

\record{Impression that answers the clinical question}

\record{Question status: answered, partly answered or inconclusive}

\record{Follow-up action, responsible clinician or team, and due date}

\record{Urgent communication: recipient, time and method}

\small\href{https://www.aaojournal.org/article/S0161-6420(19)32902-5/fulltext/}{Clinical reference for this scenario}

\normalsize

\newpage

{\small IMAGING REFERRAL AND REPORTING WORKBOOK \hfill 15 of 16}

\section*{Fundus photography}

\subsection*{What the image shows and means}

\textbf{Clinical question}\par Are there visible signs of diabetic retinopathy on screening?\par

\textbf{Microaneurysms}\par May be early diabetic retinal changes\par

\textbf{Retinal haemorrhages}\par Contribute to severity grading\par

\textbf{Abnormal new vessels}\par Suggest proliferative disease\par

\textbf{What the image leaves unanswered}\par A poor-quality photo cannot be graded reliably. Photography has limits for assessing retinal thickening.\par

\textbf{Illustration guide for a 4:3 plate}\par Show a labelled fundus view with an inset for small lesions. Label the optic disc and macula.\par

\textbf{Worked interpretation example}\par Finding: Microaneurysms. Interpretation: may be early diabetic retinal changes. If this finding is absent, use the limitations above to decide whether the original question remains open.\par

\record{Finding, location, measurement and unit}

\record{Does the finding support the hypothesis, make it less likely, or leave it unresolved?}

\record{Plain-language explanation}

\record{Audit challenge: reviewer and date, finding disputed, evidence, response or resolution}

\small\href{https://www.aaojournal.org/article/S0161-6420(19)32902-5/fulltext/}{Clinical reference for this scenario}

\normalsize

\newpage

{\small IMAGING REFERRAL AND REPORTING WORKBOOK \hfill 15 of 16}

\section*{Fundus photography}

\subsection*{Find similar cases and check the evidence}

\textbf{Clinical question}\par Are there visible signs of diabetic retinopathy on screening?\par

\textbf{Optional review aid}\par Use de-identified clinical terms. Keep names, identifiers, reports and patient images out of public searches. Record examples for review without delaying care.\par

\textbf{1. Check the decision using only pre-scan information}\par diabetic retinopathy retinal photography screening referral guideline\par

\textbf{2. Find a teaching case using the imaging findings}\par fundus photograph diabetic retinopathy microaneurysms teaching case\par

\record{Actual search terms, source link and access date}

\record{What matches and what differs}

\record{How the diagnosis was established, or whether confirmation was unclear}

\record{What the source adds and what remains uncertain}

\small\href{https://www.aaojournal.org/article/S0161-6420(19)32902-5/fulltext/}{Clinical reference for this scenario}

\normalsize

\newpage

{\small IMAGING REFERRAL AND REPORTING WORKBOOK \hfill 16 of 16}

\section*{Fluorescein angiography}

\subsection*{Before the scan}

\textbf{Clinical question}\par Which retinal vessels are leaking or failing to fill normally?\par

\textbf{What to check first}\par Review retinal exam, fundus images and OCT. Specify the remaining vascular question and whether the answer would affect care.\par

\textbf{When to proceed or defer}\par Use when timed dye imaging can clarify a meaningful question. Defer if existing examination and imaging already answer it adequately.\par

\textbf{When to defer or adapt the test}\par Review prior fluorescein reactions and pregnancy with the eye team. For a previous severe reaction, obtain specialist review before re-exposure and assess alternatives.\par

\record{Patient ID, date and ordering clinician}

\record{History received: exact account, source and relationship, recipient, date and time}

\record{Clinician assessment and prior tests, with dates and results}

\record{Guideline and exact scenario matched}

\record{Finding that supports ordering or deferral}

\record{Patient-specific reason for an exception, if applicable}

\record{Order or defer, urgency, and what the result would change}

\record{Safety decision and plain-language explanation for the patient}

\textbf{Basis for request}\par Supported clinical use, uncertain clinical benefit, research, or reason awaiting clarification. For research, record the study reference and applicable approvals.\par

\small\href{https://www.aaojournal.org/article/S0161-6420(19)32902-5/fulltext/}{Clinical reference for this scenario}

\normalsize

\newpage

{\small IMAGING REFERRAL AND REPORTING WORKBOOK \hfill 16 of 16}

\section*{Fluorescein angiography}

\subsection*{The scan-specific report}

\textbf{Clinical question}\par Which retinal vessels are leaking or failing to fill normally?\par

\textbf{What this report should record}\par Record eye, imaging phases, location and extent of leakage or nonfilling, image quality and relevant limitations.\par

\record{Patient ID, date and reporting clinician}

\record{Exam performed, technique, side or site, and comparisons}

\record{Findings and relevant measurements}

\record{Image quality, coverage or other limitations}

\record{Impression that answers the clinical question}

\record{Question status: answered, partly answered or inconclusive}

\record{Follow-up action, responsible clinician or team, and due date}

\record{Urgent communication: recipient, time and method}

\small\href{https://www.aaojournal.org/article/S0161-6420(19)32902-5/fulltext/}{Clinical reference for this scenario}

\normalsize

\newpage

{\small IMAGING REFERRAL AND REPORTING WORKBOOK \hfill 16 of 16}

\section*{Fluorescein angiography}

\subsection*{What the image shows and means}

\textbf{Clinical question}\par Which retinal vessels are leaking or failing to fill normally?\par

\textbf{Increasing spread of dye over time}\par Supports leakage\par

\textbf{Persistent capillary nonfilling}\par May indicate poor perfusion\par

\textbf{Late hyperfluorescence}\par Compare phases to distinguish leakage from staining\par

\textbf{What the image leaves unanswered}\par Blood can block the visible signal. A dark area therefore needs comparison with other images before being called nonperfusion.\par

\textbf{Illustration guide for a 4:3 plate}\par Pair early and late frames of the same area and mark the changing dye pattern.\par

\textbf{Worked interpretation example}\par Finding: Increasing spread of dye over time. Interpretation: supports leakage. If this finding is absent, use the limitations above to decide whether the original question remains open.\par

\record{Finding, location, measurement and unit}

\record{Does the finding support the hypothesis, make it less likely, or leave it unresolved?}

\record{Plain-language explanation}

\record{Audit challenge: reviewer and date, finding disputed, evidence, response or resolution}

\small\href{https://www.aaojournal.org/article/S0161-6420(19)32902-5/fulltext/}{Clinical reference for this scenario}

\normalsize

\newpage

{\small IMAGING REFERRAL AND REPORTING WORKBOOK \hfill 16 of 16}

\section*{Fluorescein angiography}

\subsection*{Find similar cases and check the evidence}

\textbf{Clinical question}\par Which retinal vessels are leaking or failing to fill normally?\par

\textbf{Optional review aid}\par Use de-identified clinical terms. Keep names, identifiers, reports and patient images out of public searches. Record examples for review without delaying care.\par

\textbf{1. Check the decision using only pre-scan information}\par suspected retinal vascular leakage fluorescein angiography indication guideline\par

\textbf{2. Find a teaching case using the imaging findings}\par fluorescein angiography retinal capillary nonperfusion late leakage teaching case\par

\record{Actual search terms, source link and access date}

\record{What matches and what differs}

\record{How the diagnosis was established, or whether confirmation was unclear}

\record{What the source adds and what remains uncertain}

\small\href{https://www.aaojournal.org/article/S0161-6420(19)32902-5/fulltext/}{Clinical reference for this scenario}

\normalsize

\end{document}
